\documentclass[a4paper, 12pt]{memoir}

\input{../packages.tex}

\title{general relativity}
\date{\today}

\newcommand{\subt}{python computations.}

\begin{document}

\frontmatter

\input{../titlepage.tex}

\tableofcontents

\mainmatter

\begin{pycode}

import sympy as sp

sp.init_printing()

class Manifold: 
    def __init__(self):
        self.g = sp.zeros(4, 4)
        self.c = [sp.zeros(4, 4),  sp.zeros(4, 4), sp.zeros(4, 4), sp.zeros(4, 4)]
        riemann1 = [sp.zeros(4, 4),  sp.zeros(4, 4), sp.zeros(4, 4), sp.zeros(4, 4)]
        riemann2 = [sp.zeros(4, 4),  sp.zeros(4, 4), sp.zeros(4, 4), sp.zeros(4, 4)]
        riemann3 = [sp.zeros(4, 4),  sp.zeros(4, 4), sp.zeros(4, 4), sp.zeros(4, 4)]
        riemann4 = [sp.zeros(4, 4),  sp.zeros(4, 4), sp.zeros(4, 4), sp.zeros(4, 4)]
        self.riemann = [riemann1, riemann2, riemann3, riemann4]
        self.ric = sp.zeros(4, 4)
        self.scal = 0
        self.ein = sp.zeros(4, 4)
        self.einup = sp.zeros(4, 4)
        self.ricup = sp.zeros(4, 4)

    def setmetric(self, g00, g11, g22, g33):
        self.g[0,0] = sp.simplify(g00)
        self.g[1,1] = sp.simplify(g11)
        self.g[2,2] = sp.simplify(g22)
        self.g[3,3] = sp.simplify(g33)

    def printmet(self):
        print(r'g =', sp.latex(self.g), '~,')
        print(r'\end{equation*}\begin{equation*}')

    def printmetinv(self):
        g_inv = self.g.inv()
        print(r'g =', sp.latex(g_inv), '~,')
        print(r'\end{equation*}\begin{equation*}')

    def setchri(self):
        g_inv = self.g.inv()
        for i in range(4):
            for j in range(4):
                for k in range(4):
                    for l in range(4):
                        self.c[k][i,j] += 1/2 * g_inv[k,l] * (sp.diff(self.g[i,l], x[j]) + sp.diff(self.g[j,l], x[i])-sp.diff(self.g[i,j], x[l]))
                        self.c[k][i,j] = sp.simplify(self.c[k][i,j])

    def printchri(self):
        count = 0
        for i in range(4):
            for j in range(4):
                for k in range(4):
                    if self.c[i][j,k] != 0:
                        print(r'\Gamma^{', i, '}_{', j, k,'} = ', sp.latex(self.c[i][j,k]), '~, \quad')
                        count += 1
                        if count % 3 == 0: 
                            print(r'\end{equation*}\begin{equation*}') 

    def setriem(self):
        g_inv = self.g.inv()
        for i in range(4):
            for j in range(4):
                for k in range(4):
                    for l in range(4):
                        self.riemann[i][j][k,l] += sp.diff(self.c[i][j,l], x[k]) - sp.diff(self.c[i][j,k], x[l])
                        for a in range(4):
                            self.riemann[i][j][k,l] += self.c[a][j,l] * self.c[i][a,k] - self.c[a][j,k] * self.c[i][a,l] 
                        self.riemann[i][j][k,l] = sp.simplify(self.riemann[i][j][k,l])

    def printriem(self):
        for i in range(4):
            for j in range(4):
                for k in range(4):
                    for l in range(4):
                        if self.riemann[i][j][k,l] != 0:
                            print(r'R^{', i, '}_{ \phantom ', i, j, k, l, '} = ', sp.latex(self.riemann[i][j][k,l]), '~,')
                            print(r'\end{equation*}\begin{equation*}')

    def setric(self): 
        for i in range(4):
            for j in range(4):
                for k in range(4):
                    self.ric[i,j] += sp.diff(self.c[k][i,j], x[k]) - sp.diff(self.c[k][i,k], x[j])
                    for a in range(4):
                        self.ric[i,j] += self.c[a][i,j] * self.c[k][a,k] - self.c[a][i,k] * self.c[k][a,j]
                    self.ric[i,j] = sp.simplify(self.ric[i,j])
    
    def setricup(self): 
        g_inv = self.g.inv()
        for i in range(4):
            for j in range(4):
                for a in range(4):
                    self.ricup[i,j] += g_inv[i,a] * self.ric[a,j]
                self.ricup[i,j] = sp.simplify(self.ricup[i,j])

    def printric(self):
        for i in range(4):
            for j in range(4):
                if self.ric[i,j] != 0:
                    print(r'R_{', i, j,'} = ', sp.latex(self.ric[i,j]), '~,')
                    print(r'\end{equation*}\begin{equation*}')

    def printricup(self):
        for i in range(4):
            for j in range(4):
                if self.ricup[i,j] != 0:
                    print(r'R^{', i, '}_{\phantom', i, j,'} = ', sp.latex(self.ricup[i,j]), '~,')
                    print(r'\end{equation*}\begin{equation*}')

    def setscal(self):
        g_inv = self.g.inv()
        for i in range(4):
            self.scal += g_inv[i,i] * self.ric[i,i]
            self.scal = sp.simplify(self.scal)

    def printscal(self):
        print(r'R = ', sp.latex(self.scal))

    def setein(self): 
        for i in range(4):
            for j in range(4):
                self.ein[i,j] += self.ric[i,j] - 0.5 * self.g[i,j] * self.scal
                self.ein[i,j] = sp.simplify(self.ein[i,j])

    def seteinup(self): 
        g_inv = self.g.inv()
        for i in range(4):
            for j in range(4):
                for a in range(4):
                    self.einup[i,j] += g_inv[i,a] * self.ein[a,j]
                self.einup[i,j] = sp.simplify(self.einup[i,j])

    def printein(self):
        for i in range(4):
            for j in range(4):
                if self.ein[i,j] != 0:
                    print(r'G_{', i, j,'} = ', sp.latex(self.ein[i,j]), '~,')
                    print(r'\end{equation*}\begin{equation*}')

    def printeinup(self):
        for i in range(4):
            for j in range(4):
                if self.einup[i,j] != 0:
                    print(r'G^{', i, '}_{\phantom',i,j,'} = ', sp.latex(self.einup[i,j]), '~,')
                    print(r'\end{equation*}\begin{equation*}')

t = sp.Symbol('t')
r = sp.Symbol('r')
theta = sp.Symbol("theta") 
phi = sp.Symbol('phi') 
x = [t,r,theta,phi]
k = sp.Symbol('k')
G = sp.Symbol('G')
M = sp.Symbol('M')

A = sp.Function('A')(r)
B = sp.Function('B')(r)
a = sp.Function('a')(t)
l = sp.Function('lambda')(r)
nu = sp.Function('nu')(r)

ES = Manifold() 
ES.setmetric(-A, B, r**2, r**2 * sp.sin(theta)**2)
ES.setchri()
ES.setriem()
ES.setric()
ES.setricup()
ES.setscal()
ES.setein()
ES.seteinup()

IS = Manifold() 
IS.setmetric(- sp.exp(nu), sp.exp(l), r**2, r**2 * sp.sin(theta)**2)
IS.setchri()
IS.setriem()
IS.setric()
IS.setricup()
IS.setscal()
IS.setein()
IS.seteinup()

FRW = Manifold() 
FRW.setmetric(-1, a**2 / (1 - k * r**2), a**2 * r**2, a**2 * r**2 * sp.sin(theta)**2)
FRW.setchri()
FRW.setriem()
FRW.setric()
FRW.setricup()
FRW.setscal()
FRW.setein()
FRW.seteinup()

def plot3(x, f, g, h, rangex, rangey, fig, leg, negx, negy):
    rangexx = rangex
    rangeyy = rangey
    if negx == True:
        rangexx = 0
    if negy == True:
        rangeyy = 0
    x = sp.Symbol('x')
    p = sp.plot((f, (x, -rangexx, rangex)), (g, (x, -rangexx, rangex)), (h, (x, -rangexx, rangex)), ylim=[-rangeyy, rangey], legend= leg, show=False, line_color='#00AFB3')
    p[0].line_color='red'
    p[1].line_color='violet'
    p[2].line_color='blue'
    p.save(f'fig/fig{fig}.pgf')
    print(r'\input{fig/fig'+ rf'{fig}' + r'.pgf}')
    
def plot2(x, f, g, rangex, rangey, fig, leg, negx, negy):
    rangexx = rangex
    rangeyy = rangey
    if negx == True:
        rangexx = 0
    if negy == True:
        rangeyy = 0
    x = sp.Symbol('x')
    p = sp.plot((f, (x, -rangexx, rangex)), (g, (x, -rangexx, rangex)), ylim=[-rangeyy, rangey], legend= leg, show=False, line_color='#00AFB3')
    p[0].line_color='red'
    p[1].line_color='blue'
    p.save(f'fig/fig{fig}.pgf')
    print(r'\input{fig/fig'+ rf'{fig}' + r'.pgf}')

\end{pycode} 



\chapter*{Abstract}

    In these notes, we will compute general relativity with the help of python. We will take the exterior Schwarschild metric, the interior Schwarschild metric and the Friedmann-Robertson-Walker metric. We will calculate the Christoffel symbols, the Riemann tensor, the Ricci tensor and the Einstein tensor.

\chapter{Exterior Schwarschild metric}

    The metric is 
    \begin{equation*}
        \pyc{ES.printmet()} ~.
    \end{equation*}

    Its inverse is 
    \begin{equation*}
        \pyc{ES.printmetinv()} ~.
    \end{equation*}

\section{Christoffel symbols}
    The non-vanishing components of the Christoffel symbols are
    \begin{equation*}
        \pyc{ES.printchri()} ~.
    \end{equation*}

\section{Riemann tensor}
    The non-vanishing components of the Riemann tensor are
    \begin{equation*}
        \pyc{ES.printriem()} ~.
    \end{equation*}

\section{$(0,2)$ Ricci tensor}
    The non-vanishing components of the $(0,2)$ Ricci tensor are
    \begin{equation*}
        \pyc{ES.printric()} ~.
    \end{equation*}

\section{$(1,1)$ Ricci tensor}
    The non-vanishing components of the $(1,1)$ Ricci tensor are
    \begin{equation*}
        \pyc{ES.printricup()} ~.
    \end{equation*}

\section{Ricci scalar}
    The Ricci scalar is
    \begin{equation*}
        \pyc{ES.printscal()} ~.
    \end{equation*}

\section{$(0,2)$ Einstein tensor}
    The non-vanishing components of the $(0,2)$ Einstein tensor are
    \begin{equation*}
        \pyc{ES.printein()} ~.
    \end{equation*}

\section{$(1,1)$ Einstein tensor}
    The non-vanishing components of the $(1,1)$ Einstein tensor are
    \begin{equation*}
        \pyc{ES.printeinup()} ~.
    \end{equation*}

\chapter{Interior Schwarschild metric}

    The metric is 
    \begin{equation*}
        \pyc{IS.printmet()} ~.
    \end{equation*}

    Its inverse is 
    \begin{equation*}
        \pyc{IS.printmetinv()} ~.
    \end{equation*}

\section{Christoffel symbols}
    The non-vanishing components of the Christoffel symbols are
    \begin{equation*}
        \pyc{IS.printchri()} ~.
    \end{equation*}

\section{Riemann tensor}
    The non-vanishing components of the Riemann tensor are
    \begin{equation*}
        \pyc{IS.printriem()} ~.
    \end{equation*}

\section{$(0,2)$ Ricci tensor}
    The non-vanishing components of the $(0,2)$ Ricci tensor are
    \begin{equation*}
        \pyc{IS.printric()} ~.
    \end{equation*}

\section{$(1,1)$ Ricci tensor}
    The non-vanishing components of the $(1,1)$ Ricci tensor are
    \begin{equation*}
        \pyc{IS.printricup()} ~.
    \end{equation*}

\section{Ricci scalar}
    The Ricci scalar is
    \begin{equation*}
        \pyc{IS.printscal()} ~.
    \end{equation*}

\section{$(0,2)$ Einstein tensor}
    The non-vanishing components of the $(0,2)$ Einstein tensor are
    \begin{equation*}
        \pyc{IS.printein()} ~.
    \end{equation*}

\section{$(1,1)$ Einstein tensor}
    The non-vanishing components of the $(1,1)$ Einstein tensor are
    \begin{equation*}
        \pyc{IS.printeinup()} ~.
    \end{equation*}
    
\chapter{FRW metric}

    The metric is 
    \begin{equation*}
        \pyc{FRW.printmet()} ~.
    \end{equation*}

    Its inverse is 
    \begin{equation*}
        \pyc{FRW.printmetinv()} ~.
    \end{equation*}
    
\section{Christoffel symbols}
    The non-vanishing components of the Christoffel symbols are
    \begin{equation*}
        \pyc{FRW.printchri()} ~.
    \end{equation*}

\section{Riemann tensor}
    The non-vanishing components of the Riemann tensor are
    \begin{equation*}
        \pyc{FRW.printriem()} ~.
    \end{equation*}

\section{$(0,2)$ Ricci tensor}
    The non-vanishing components of the $(0,2)$ Ricci tensor are
    \begin{equation*}
        \pyc{FRW.printric()} ~.
    \end{equation*}

\section{$(1,1)$ Ricci tensor}
    The non-vanishing components of the $(1,1)$ Ricci tensor are
    \begin{equation*}
        \pyc{FRW.printricup()} ~.
    \end{equation*}

\section{Ricci scalar}
    The Ricci scalar is
    \begin{equation*}
        \pyc{FRW.printscal()} ~.
    \end{equation*}

\section{$(0,2)$ Einstein tensor}
    The non-vanishing components of the $(0,2)$ Einstein tensor are
    \begin{equation*}
        \pyc{FRW.printein()} ~.
    \end{equation*}

\section{$(1,1)$ Einstein tensor}
    The non-vanishing components of the $(1,1)$ Einstein tensor are
    \begin{equation*}
        \pyc{FRW.printeinup()} ~.
    \end{equation*}

\nocite{python}

\clearpage
\phantomsection
\printbibliography

\end{document}
